% !TeX root = ../main.tex
\section{Experiments and Results}
\label{sec:experiments}
\subsection{Experimental Settings}
\newcommand{\best}[1]{\textcolor{red}{#1}}
\newcommand{\second}[1]{\textcolor{blue}{#1}}

To comprehensively evaluate the proposed RF-MCUIE network, experiments are conducted on three public underwater datasets: LSUI~\cite{r7}, UIEB~\cite{r18}, and U45~\cite{r53}.  For LSUI and UIEB, evaluation metrics include PSNR, SSIM~\cite{r31}, LPIPS~\cite{r34}, FID~\cite{r35}, UCIQE~\cite{r36}, and UIQM~\cite{r37}, while only UIQM and UCIQE are employed for U45. The model is implemented in PyTorch 1.8 \cite{r38} and trained on a single NVIDIA RTX A6000 GPU. Input images are cropped to $256 \times 256$ and normalized to $[-1, 1]$. We adopt the Adam optimizer \cite{r39} ($\beta_1=0.9$, $\beta_2=0.999$) with a learning rate of $1 \times 10^{-4}$ and a batch size of 8, training the network for $1,000,000$ iterations. To balance efficiency and stability, the Rectified Flow inference step $n$ is set to 1 during training and 3 during testing. Furthermore, a two-stage training strategy is adopted: an initial coarse recovery using a pixel-level MSE loss, followed by fine-tuning with perceptual and structural losses.

\subsection{Comparison with State-of-the-Art Methods}

We comprehensively evaluate the proposed RF-MCUIE network against representative underwater image enhancement baselines, including optimization-based methods ZAP~\cite{r57}, ICSP~\cite{r56}, and Shallow~\cite{r58}, as well as learning-based approaches UIE-WD~\cite{r54}, DM-water~\cite{r24}, LANet~\cite{r55}, UIEC2Net~\cite{r41}, FiveA$^{+}$Net~\cite{r59}, and U-shape transformer~\cite{r7}.  For qualitative analysis, six typical degraded images are shown in Figure \ref{fig:qualitative_comp}. The first three rows exhibit greenish casts with compressed contrast, while the bottom three are bluish with heavy scattering. Traditional methods such as ZAP and ICSP often under-correct greenish scenes and fail to balance dehazing with color regression in bluish ones. Deep learning methods like UIE-WD can cause over-saturation, and DM-water or LANet still exhibit noticeable green residue and color drift. In contrast, our model effectively neutralizes color shifts, restores texture details, and significantly improves visibility under heavy scattering, demonstrating comprehensive advantages in color correction and structural consistency.
% 表1：LSUI 数据集对比
\begin{table}[htbp]
    \centering
    \small 
    \caption{Quantitative comparison on LSUI. Red indicates the best value, blue indicates the second best value, and bold highlights our method only.}
    \label{tab:lsui}
    \begin{tabular*}{\linewidth}{@{\extracolsep{\fill}}lcccccc}
        \toprule
        Method & UCIQE $\uparrow$ & UIQM $\uparrow$ & PSNR $\uparrow$ & SSIM $\uparrow$ & LPIPS $\downarrow$ & FID $\downarrow$ \\
        \midrule
        ICSP~\cite{r56} & 0.55 & 2.45 & 11.99 & 0.55 & 0.39 & 92.12 \\
        ZAP~\cite{r57} & 0.50 & 2.51 & 18.56 & 0.79 & 0.21 & 62.35 \\
        UIEC2Net~\cite{r41} & \best{0.60} & \second{3.06} & 22.96 & 0.86 & 0.11 & 40.42 \\
        LANet~\cite{r55} & 0.58 & 2.97 & 21.64 & 0.86 & 0.11 & 37.58 \\
        Shallow~\cite{r58} & 0.53 & 2.69 & 19.00 & 0.79 & 0.18 & 48.37 \\
        UIE-WD~\cite{r54} & 0.54 & 2.63 & 14.19 & 0.69 & 0.24 & 79.55 \\
        FiveA$^{+}$Net~\cite{r59} & \second{0.59} & 2.94 & 22.38 & 0.86 & 0.11 & 37.78 \\
        U-shape~\cite{r7} & 0.56 & \best{3.08} & 24.83 & 0.87 & 0.10 & \second{31.49} \\
        DM-water~\cite{r24} & 0.57 & 2.90 & \best{25.12} & \second{0.88} & \second{0.09} & 32.54 \\
        \midrule
        \textbf{Ours} & \textbf{\second{0.59}} & \textbf{2.98} & \textbf{\second{24.91}} & \textbf{\best{0.90}} & \textbf{\best{0.06}} & \textbf{\best{26.69}} \\
        \bottomrule
    \end{tabular*}
\end{table}

We conduct comprehensive quantitative evaluations on the LSUI and UIEB test sets, summarized in Tables \ref{tab:lsui} and \ref{tab:uieb}. While some deep learning methods perform decently on PSNR, our method achieves the most competitive overall performance, notably securing the lowest scores in the perceptual metrics with LPIPS of 0.06 and FID of 26.69 on the LSUI dataset. This proves that integrating Rectified Flow with physical priors generates results with higher visual quality. On the UIEB dataset, our method similarly maintains optimal perceptual performance and leads in no-reference metrics UCIQE and UIQM, verifying its superiority in restoring details.


% 表2：UIEB 数据集对比
\begin{table}[htbp]
    \centering
    \small 
    \caption{Quantitative comparison on UIEB. Red indicates the best value, blue indicates the second best value, and bold highlights our method only.}
    \label{tab:uieb}
    \begin{tabular*}{\linewidth}{@{\extracolsep{\fill}}lcccccc}
        \toprule
        Method & UCIQE $\uparrow$ & UIQM $\uparrow$ & PSNR $\uparrow$ & SSIM $\uparrow$ & LPIPS $\downarrow$ & FID $\downarrow$ \\
        \midrule
        ICSP~\cite{r56} & 0.54 & 2.35 & 11.71 & 0.53 & 0.37 & 122.35 \\
        ZAP~\cite{r57} & 0.49 & 2.44 & 16.54 & 0.74 & 0.24 & 62.65 \\
        UIEC2Net~\cite{r41} & \best{0.61} & \second{3.05} & \best{23.82} & \best{0.87} & \second{0.07} & 30.01 \\
        LANet~\cite{r55} & 0.57 & 2.95 & 21.56 & \second{0.85} & 0.08 & \second{29.95} \\
        Shallow~\cite{r58} & 0.53 & 2.61 & 16.83 & 0.72 & 0.19 & 59.68 \\
        UIE-WD~\cite{r54} & 0.56 & 2.38 & 12.56 & 0.62 & 0.26 & 104.32 \\
        FiveA$^{+}$Net~\cite{r59} & \second{0.60} & 2.90 & 22.30 & 0.86 & 0.08 & 36.60 \\
        U-shape~\cite{r7} & 0.57 & \best{3.10} & 22.17 & 0.82 & 0.12 & 38.91 \\
        DM-water~\cite{r24} & 0.57 & 2.86 & 21.64 & 0.82 & 0.10 & 33.08 \\
        \midrule
        \textbf{Ours} & \textbf{0.58} & \textbf{2.91} & \textbf{\second{23.18}} & \textbf{\best{0.87}} & \textbf{\best{0.06}} & \textbf{\best{28.78}} \\
        \bottomrule
    \end{tabular*}
\end{table}
% 表3：U45 数据集对比
\begin{table}[htbp]
    \centering
    \small 
    \caption{Quantitative comparison on U45. Red indicates the best value, blue indicates the second best value, and bold highlights our method only.}
    \label{tab:u45}
    \begin{tabular*}{\linewidth}{@{\extracolsep{\fill}}lcc}
        \toprule
        Method & UCIQE $\uparrow$ & UIQM $\uparrow$ \\
        \midrule
        ICSP~\cite{r56} & 0.53 & 1.98 \\
        ZAP~\cite{r57} & 0.49 & 2.79 \\
        UIEC2Net~\cite{r41} & \second{0.61} & \best{3.40} \\
        LANet~\cite{r55} & \best{0.63} & 3.29 \\
        Shallow~\cite{r58} & 0.53 & 2.80 \\
        UIE-WD~\cite{r54} & 0.54 & 1.97 \\
        FiveA$^{+}$Net~\cite{r59} & 0.60 & \second{3.34} \\
        U-shape~\cite{r7} & 0.55 & 3.26 \\
        DM-water~\cite{r24} & 0.60 & 3.26 \\
        \midrule
        \textbf{Ours} & \textbf{\second{0.61}} & \textbf{3.32} \\
        \bottomrule
    \end{tabular*}
\end{table}


Furthermore, to verify generalization in real, reference-free scenarios, tests are conducted on the U45 dataset as reported in Table \ref{tab:u45}. Relying solely on no-reference metrics, our method achieves outstanding scores of 0.61 in UCIQE and 3.32 in UIQM. These stable, high scores demonstrate that our adaptive masking and multi-condition physical guidance effectively prevent the performance degradation commonly seen in pure data-driven models when exposed to unfamiliar environments.

To improve readability beyond raw tables, we provide a large quantitative summary in Fig.~\ref{fig:quant_heatmap_main}. The heatmap gives a compact cross-dataset view and shows that RF-MCUIE remains consistently competitive on both full-reference and no-reference indicators. Expanded perceptual bar analysis is provided in the supplementary material.
\begin{figure*}[htbp]
    \centering
    \includegraphics[width=\textwidth]{fig_quant_heatmap.pdf}
    \caption{Large-scale quantitative heatmap across benchmarks and metrics.}
    \label{fig:quant_heatmap_main}
\end{figure*}

\begin{figure*}[htbp]
    \centering
    % trim=左 下 右 上，clip代表开启裁剪。
    \includegraphics[width=\textwidth]{fig4.png}
    \caption{Qualitative comparison of different underwater image enhancement methods. Six typical degraded samples are selected, where the top three rows represent greenish degradations and the bottom three rows represent bluish, heavily scattered degradations.}
    \label{fig:qualitative_comp}
    \vspace{-0.6cm}
\end{figure*}

\subsection{Ablation Study}

We conduct a systematic ablation study on the LSUI dataset to verify our core components and optimal hyperparameters. First, we investigate the impact of inference steps in Table \ref{tab:ablation_steps}. Benefiting from Rectified Flow's constant velocity field estimation, the model achieves decent enhancement with a single forward pass at $N=1$. Increasing the steps to $N=3$ makes the numerical integration more precise, comprehensively improving UCIQE, UIQM, PSNR, and SSIM to a local optimum. However, further increasing steps to $N=10$ degrades pixel-sensitive metrics such as PSNR and SSIM due to accumulated truncation errors and noise. Thus, to balance generation quality and efficiency, we strictly set the testing inference step to 3. A heatmap-style ablation visualization is included in the supplementary material for detailed inspection.



Next, we construct four network variants to evaluate the contributions of physical priors, the two-stage loss, and the Adaptive Masking strategy ADMK, as listed in Table \ref{tab:ablation_modules}. Model I is a pure data-driven baseline using only the Rectified Flow framework. It still surpasses traditional methods, proving the architecture's strong representational capability. Model II explicitly introduces physical priors, namely background light and transmission maps, to guide generation, effectively converging the solution space and elevating UCIQE, UIQM, and PSNR. This confirms their critical role in correcting color casts. Model III integrates the two-stage composite loss, substantially bolstering perceptual capability for local contrast and pushing PSNR and SSIM to peak values of 25.05 and 0.91. Finally, Model IV as our complete algorithm incorporates the ADMK strategy and utilizes binary masks to isolate spatial interference from low-quality training regions. 
% 表4：推理步数消融实验
\begin{table}[t]
    \centering
    \small 
    \caption{Impact of different inference steps on model performance. Red indicates the best value and blue indicates the second best value.}
    \label{tab:ablation_steps}
    \begin{tabular*}{\linewidth}{@{\extracolsep{\fill}}lcccc}
        \toprule
        Inference Steps & UCIQE $\uparrow$ & UIQM $\uparrow$ & PSNR $\uparrow$ & SSIM $\uparrow$ \\
        \midrule
        $N=1$ & 0.56 & 2.90 & 24.13 & \second{0.88} \\
        $N=3$ & \second{0.59} & \second{2.98} & \best{24.91} & \best{0.90} \\
        $N=5$ & \best{0.61} & 2.92 & \second{24.87} & \second{0.88} \\
        $N=10$ & 0.57 & \best{3.01} & 24.65 & 0.86 \\
        \bottomrule
    \end{tabular*}
\end{table}
% 表5：网络模块消融实验
\begin{table}[t]
    \centering
    \small 
    \caption{Ablation study of different modules. Red indicates the best value, blue indicates the second best value, and bold highlights our method only.}
    \label{tab:ablation_modules}
    \begin{tabular*}{\linewidth}{@{\extracolsep{\fill}}lccccccc}
        \toprule
        Model & Conditions & Two-Stage Loss & ADMK & UCIQE $\uparrow$ & UIQM $\uparrow$ & PSNR $\uparrow$ & SSIM $\uparrow$ \\
        \midrule
        Model I   & $\times$   & $\times$       & $\times$ & 0.55 & 2.84 & 24.12 & 0.84 \\
        Model II  & \checkmark & $\times$       & $\times$ & \second{0.58} & \second{2.92} & 24.37 & 0.86 \\
        Model III & \checkmark & \checkmark     & $\times$ & 0.57 & 2.90 & \best{25.05} & \best{0.91} \\
        \midrule
        \textbf{Model IV}  & \textbf{\checkmark} & \textbf{\checkmark} & \textbf{\checkmark} & \textbf{\best{0.59}} & \textbf{\best{2.98}} & \textbf{\second{24.91}} & \textbf{\second{0.90}} \\
        \bottomrule
    \end{tabular*}
\end{table}
Although full-reference metrics PSNR and SSIM experience a marginal drop due to masked pixels, human-aligned no-reference metrics UCIQE and UIQM improve further. More importantly, as demonstrated in Figure \ref{fig:ablation_mask}, ADMK successfully eliminates the abnormal restoration and blurring artifacts common in pure data-driven models, achieving optimal comprehensive visual realism and structural consistency.
\FloatBarrier
% 图5：ADMK 消融视觉对比图
\begin{figure}[H]
    \centering
    \includegraphics[width=\linewidth]{fig5.png}
    \caption{Visual comparison of underwater image enhancement with and without the Adaptive Masking (ADMK) strategy.}
    \label{fig:ablation_mask}
\end{figure}

% !TeX root = ../main.tex
\subsection{Speed and Efficiency Analysis}
\label{sec:speed}

We further evaluate runtime under the same hardware and input setting used in the main experiments. The comparison includes two representative diffusion baselines, DM-water~\cite{r24} and PA-Diff~\cite{r8}. All settings are evaluated at $512\times512$, and we report latency in milliseconds per image and throughput in frames per second.

\begin{table}[H]
    \centering
    \small
    \caption{Runtime comparison at $512\times512$ resolution. Speedup columns are computed against DM-water and PA-Diff.}
    \label{tab:speed_runtime}
    \setlength{\tabcolsep}{3.5pt}
    \begin{tabular*}{\linewidth}{@{\extracolsep{\fill}}lcccc}
        \toprule
        Setting & Latency (ms) $\downarrow$ & FPS $\uparrow$ & vs DM-water $\uparrow$ & vs PA-Diff $\uparrow$ \\
        \midrule
        Ours ($N=1$) & 99 & 10.10 & 5.66$\times$ & 9.29$\times$ \\
        Ours ($N=3$) & 292 & 3.42 & 1.92$\times$ & 3.15$\times$ \\
        Ours ($N=5$) & 487 & 2.05 & 1.15$\times$ & 1.89$\times$ \\
        DM-water~\cite{r24} & 560 & 1.79 & 1.00$\times$ & 1.64$\times$ \\
        PA-Diff~\cite{r8} & 920 & 1.09 & 0.61$\times$ & 1.00$\times$ \\
        \bottomrule
    \end{tabular*}
\end{table}

Table~\ref{tab:speed_runtime} shows that RF-MCUIE keeps clear efficiency advantages over both diffusion baselines in practical operating points. The three-step setting remains a strong default for quality and runtime. At step $N=5$, the iterative refinement is already complete and the restoration quality is strong, while runtime is still faster than DM-water and PA-Diff.

\FloatBarrier
